\section{Introduction}
3D Indoor Scene Generation is a foundational task in computer graphics, virtual reality, and synthetic data generation for embodied AI. Since 2024, diffusion models have emerged as the leading paradigm for this task. State-of-the-art models like CommonScenes, \cite{zhai2023commonscenesgeneratingcommonsense3d} EchoScene, \cite{zhai2025echosceneindoorscenegeneration} and DiffuScene \cite{tang2024diffuscenedenoisingdiffusionmodels} leverage continuous diffusion processes to arrange 3D objects realistically based on learned spatial distributions, outputting a set of oriented bounding boxes (OBBs) that dictate the position, scale, and orientation of furniture objects in a 3D scene.

\textbf{Challenges:} Despite their visual appeal and strong semantic layout understanding, a persistent challenge with unconstrained 3D diffusion models is the lack of strict physical guarantees. Because the diffusion training objective typically minimizes a reconstruction loss over bounding box parameters, it does not explicitly penalize inter-object collisions or out-of-room placements. Existing methods that aim to address these inaccuracies typically involve retraining the network entirely with differentiable collision constraints. This approach is computationally expensive, requires access to model weights, and can destabilize the training process.

While recent work such as PhyScene \cite{yang2024physcenephysicallyinteractable3d} has attempted to improve physical plausibility by incorporating collision-aware objectives into diffusion model training or by applying gradient-based guidance during inference, these approaches introduce several limitations. First, modifying the training objective can interfere with the model's ability to learn the underlying semantic distribution of indoor layouts, potentially reducing generative diversity. Second, retraining or fine-tuning diffusion models is computationally expensive and often impractical, particularly for large-scale pretrained models, and any such modification is tied to a single architecture, offering no guarantee of transfer to other generative frameworks. 

\textbf{Our Approach:} To address these limitations, we propose a \textit{Collision Resolver}, a lightweight and fully training-free framework for enforcing physical plausibility in generated 3D scene layouts. Unlike approaches that incorporate physical constraints during training or the diffusion process, our method operates exclusively as a post-processing step on the final denoised layout, directly resolving inter-object collisions and invalid furniture placements. Operating entirely at inference time, our method requires no retraining and therefore preserves the generative model's learned distribution.
This strict separation from the generation process makes our approach inherently model-agnostic and enables it to be readily applied to layouts produced by different 3D scene generation architectures.

We instantiate and validate this idea using DiffuScene~\cite{tang2024diffuscenedenoisingdiffusionmodels} as our primary baseline, applying a deterministic, alignment-biased SAT-based Collision Resolver directly to its generated bounding boxes. Crucially, because the resolver operates purely on layout geometry - object extents, positions, and orientations - rather than on any pipeline-specific representation, the same module can be dropped into the output stage of other layout-generating scene synthesis methods with no retraining. This model-agnostic formulation is central to our contribution: rather than proposing a collision-resolution heuristic tuned to one model's failure modes, we show that treating collision resolution as a purely geometric, post-hoc problem yields a solver that is reusable across architecturally distinct scene synthesis methods.

Concretely, this paper makes the following contributions:
\begin{itemize}
    \item We formulate collision resolution for generated indoor scenes as a deterministic, training-free post-processing problem, decoupling physical plausibility entirely from the generative model's learning objective.
    \item We introduce an alignment-biased extension to the classical Separating Axis Theorem that jointly enforces room-boundary containment and resolves pairwise OBB collisions. This ultimately removes inter-object collisions and minimizes unnatural displacement of furniture.
    \item We demonstrate the effectiveness of this resolver on the DiffuScene pipeline, achieving substantial reductions in object- and scene-level collision rates while preserving generative quality without any retraining.
\end{itemize}